\documentclass[11pt]{article}
\usepackage[margin=1in]{geometry}
\usepackage{amsmath,amssymb}
\usepackage{booktabs}
\usepackage{graphicx}
\usepackage{hyperref}
\usepackage{url}
\usepackage{authblk}

\title{Benchmark Verdicts for Analog Ising Machines Invert Under Tuning:\\
A Negative Result and a Reproducible Protocol}

\author[1]{Shmuel Y.\ Biller}
\affil[1]{Independent researcher}

\date{\today}

\begin{document}
\maketitle

\begin{abstract}
We report a systematic simulation study of a measurement-feedback parametron
Ising machine (an ``analog'' or coherent-Ising-machine-class solver) benchmarked
against discrete Simulated Bifurcation (dSBM). Our central finding is
methodological and negative: on identical instances at matched compute, the
competitive verdict \emph{inverted twice} depending solely on which solver
received parameter tuning. Untuned analog versus untuned digital gave dSBM a
$5.9$--$7.4\times$ advantage; tuned analog versus untuned digital gave the
analog machine a $2.3$--$6.8\times$ advantage; tuning both --- the only
defensible comparison --- gave near-parity, and we further find that the
\emph{sign} of the residual difference depends on how large a tuning grid each
side receives. We argue that untuned baselines are a systemic hazard in
this literature and propose a minimal protocol to avoid them. We additionally
report: (i) a measured time-to-solution scaling of $\mathrm{TTS}\sim N^{2.115}$
($95\%$ CI $[1.93,2.30]$, $\Delta\mathrm{AIC}=9.97$ favouring a power law over
exponential growth); (ii) four algorithmic modifications yielding a compounded
$\approx\!24\times$ improvement over a schedule frozen at small $N$, of which the
largest arises from a parameter never previously varied; and (iii) an analysis
showing that the architecture's operation-count advantage is negligible in
measurement-feedback form, because the coupling multiply is digital in both
machines. All simulation code, the hashed benchmark instance library, and the
full result set are released.
\end{abstract}

\section{Introduction}

Analog and physics-inspired Ising machines are widely proposed as accelerators
for combinatorial optimization. Coherent Ising machines with measurement
feedback have been demonstrated at large spin counts at room temperature, and a
substantial literature reports favourable comparisons against digital solvers.

A recurring weakness in such comparisons is that the proposed machine is
typically tuned --- often extensively --- while the digital baseline is run at
default or literature settings. We encountered this failure mode directly and
severely: over the course of this study our own competitive conclusion reversed
twice, purely as a function of which side had been tuned, before settling at
near-parity. We report this as the primary result, because we believe it is
common and largely unreported.

We study a parametron (degenerate parametric oscillator) network with digital
measurement feedback, the architecture of Goto's parametron \cite{goto1959}
combined with the measurement-feedback scheme of modern coherent Ising machines
\cite{inagaki2016,honjo2021} and Leleu-type amplitude error correction
\cite{leleu2019}. Our digital comparator is discrete Simulated Bifurcation
\cite{goto2019,goto2021}.

Our contributions:
\begin{enumerate}
\item A demonstration that benchmark verdicts in this class invert under tuning,
      with the full sequence documented (Sec.~\ref{sec:tuning}).
\item A properly controlled near-parity result (Sec.~\ref{sec:parity}).
\item A measured scaling law with confidence intervals and model comparison
      (Sec.~\ref{sec:scaling}).
\item Four algorithmic findings, including one apparently novel method
      (Sec.~\ref{sec:algo}).
\item An operation-count analysis identifying where, if anywhere, an advantage
      can exist for this architecture (Sec.~\ref{sec:where}).
\end{enumerate}

We emphasise at the outset that all results here are simulation. No hardware was
built. Claims about physical implementations are explicitly labelled as
projections in Sec.~\ref{sec:limits}.

\section{Methods}

\subsection{Model}
Each spin $i$ is represented by a slow amplitude envelope $x_i$ obeying a
rotating-frame stochastic differential equation, integrated by Euler--Maruyama:
\begin{align}
dx_i &= \left[(p(t)-1-x_i^2)x_i + \epsilon f_i\right]dt + \sigma\,dW_i,\\
f_i  &= e_i \sum_j \tilde J_{ij}\,x_j(t-\tau_{\text{lag}}),\\
de_i &= -\beta e_i (x_i^2 - a(t))\,dt, \qquad a(t)=\max(p(t)-1,\,0.05),
\end{align}
where $p(t)$ is the pump ramp, $\tilde J$ is row-sum normalised, $e_i$ are
Leleu-style amplitude error-correction variables, and the feedback uses a
measurement delayed by one update interval. Spins are read out as
$s_i=\mathrm{sign}(x_i)$.

\subsection{Comparator}
Discrete Simulated Bifurcation \cite{goto2021} with inelastic walls
($|x|\le 1$, momentum reset on contact) and $\mathrm{sign}(x)$ coupling, using a
standard coupling normalisation.

\subsection{Cost metric}
We count \emph{matrix--vector products} (equivalently, multiply--accumulate
operations), not wall-clock time or integration steps. This is the fair unit
because it is the dominant work in both algorithms and, critically, is performed
digitally in both: the analog machine's coupling field is computed on an FPGA and
dSBM's on a GPU. Time-to-solution is
$\mathrm{TTS}_{99}=C\ln(0.01)/\ln(1-p)$, with $C$ the cost of one anneal and $p$
the per-anneal success probability.

\subsection{Instances and reference}
Dense random $\pm1$ Ising instances, $J_{ij}\in\{-1,+1\}$, $J_{ii}=0$. For
$N\le22$ ground states were obtained by exhaustive enumeration; above that the
reference is the best energy found across all runs of that instance. We verified
this reference is not loose: on $80$ instances at $N=12,16,20,22$, the
best-found energy equalled the exact ground state in every case (inflation factor
$1.00$).

\subsection{Statistics}
Unless stated otherwise, $25$ instances per problem size, medians across
instances, and bootstrap confidence intervals resampled over instances. Success
probabilities are computed over $100$--$300$ independent anneals per instance.

\section{Benchmark verdicts invert under tuning}
\label{sec:tuning}

Table~\ref{tab:invert} shows the same comparison, on the same instances, at
matched matrix-multiply budget, evaluated three times during this study.

\begin{table}[h]
\centering
\begin{tabular}{lccl}
\toprule
Comparison & $N=32$ & $N=64$ & Status \\
\midrule
Untuned analog vs.\ untuned dSBM & dSBM $5.9\times$ & dSBM $7.4\times$ & analog untuned \\
Tuned analog vs.\ untuned dSBM   & analog $6.8\times$ & analog $2.3\times$ & dSBM untuned \\
\textbf{Both tuned}              & \multicolumn{2}{c}{\textbf{near-parity (see Sec.~4)}} & \textbf{defensible} \\
\bottomrule
\end{tabular}
\caption{The competitive verdict inverts twice depending only on which solver is
tuned. Only the third row is a defensible comparison. Both solvers turned out to
be running substantially longer than optimal in the first two rows.}
\label{tab:invert}
\end{table}

The magnitude matters: the first two rows differ by a factor of roughly $40$ at
$N=32$ and would support opposite conclusions in a paper abstract. Neither was
produced dishonestly; each was the result of tuning one side and not the other,
which is the default outcome when a research group benchmarks its own system
against a literature baseline.

\paragraph{Proposed minimal protocol.}
We adopt, and suggest, the following rule: \emph{no comparison is reportable
unless the same tuning budget has been applied to every algorithm in it, and the
settings used for each are stated alongside the result, together with the size
of the parameter grid searched for each.} An untuned baseline is not a
baseline, and a tuned baseline whose search budget is unstated is only
marginally better (Sec.~\ref{sec:parity}). We further recommend counting cost in operations rather than
wall-clock time, since wall-clock comparisons across substrates conflate
algorithmic quality with implementation effort.

\section{Near-parity result}
\label{sec:parity}

With run length optimised for both solvers, the two methods are at near-parity.
We deliberately do not report a single ratio here, because the residual
difference proved smaller than the effect of the tuning grid itself: two
independent tuning efforts on the same instance family gave dSBM ahead by
$1.13\times$ at $N=32$ in one case, and the analog machine ahead by $1.72\times$
in the other, with an exact tie at $N=64$. Both efforts searched comparable but
not identical grids.

We treat this as a second-order instance of the paper's main finding rather
than a contradiction of it: if the sign of a competitive verdict can be changed
by the size of a tuning grid, a single reported ratio from a single tuning
effort carries little information. The defensible claim is near-parity, and any
stronger claim in either direction requires the tuning budget searched --- not
only the final settings --- to be reported alongside the result. Fitted scaling
exponents under equal tuning are $N^{2.72}$ (analog) versus $N^{2.32}$ (dSBM),
so the deficit widens slowly with problem size.

We regard near-parity against a mature, corporate-refined production algorithm as
a substantive result for an architecture that does not physically exist. We also
regard it as insufficient to support an energy claim, for reasons given in
Sec.~\ref{sec:where}.

\section{Scaling}
\label{sec:scaling}

Time-to-solution was measured at six sizes ($N=8$--$64$), $25$ instances each,
with anneal length optimised independently at every size
(Table~\ref{tab:scaling}).

\begin{table}[h]
\centering
\begin{tabular}{rrrrr}
\toprule
$N$ & Optimal anneal & $p$ at optimum & TTS (matmuls) & Local slope \\
\midrule
8  & 3{,}000  & 0.97 & 394    & --- \\
16 & 3{,}000  & 0.51 & 1{,}937  & 2.30 \\
24 & 10{,}000 & 0.75 & 3{,}322  & 1.33 \\
32 & 10{,}000 & 0.52 & 6{,}274  & 2.21 \\
48 & 10{,}000 & 0.27 & 14{,}633 & 2.09 \\
64 & 10{,}000 & 0.11 & 39{,}518 & 3.45 \\
\bottomrule
\end{tabular}
\caption{Scaling scan, $25$ instances per size.}
\label{tab:scaling}
\end{table}

A power law fits substantially better than an exponential:
$\mathrm{TTS}\sim N^{2.115}$ (bootstrap $95\%$ CI $[1.93,2.30]$,
$\mathrm{AIC}=-18.21$) versus $\mathrm{TTS}\sim\exp(0.0748N)$
($\mathrm{AIC}=-8.24$). With $\Delta\mathrm{AIC}=9.97$ the exponential model has
relative likelihood $\approx 0.007$. Growth over the measured range is therefore
polynomial.

Two caveats. First, the local slope in the final interval ($48\to64$) is $3.45$,
the steepest in the scan; this is consistent either with noise at $25$ instances
or with genuine curvature, and data ending at $N=64$ cannot distinguish them.
Second, exponent estimates converged only as instance count rose
($2.48\to2.11\to2.01\to2.115$ at $3,5,10,25$ instances), with the confidence
interval halving between $10$ and $25$. We recommend $\ge 25$ instances per size
for scaling claims in this class.

\section{Algorithmic findings}
\label{sec:algo}

Four parameters that had been fixed by inheritance rather than measurement
yielded substantial gains (Table~\ref{tab:algo}).

\begin{table}[h]
\centering
\begin{tabular}{lll}
\toprule
Modification & Gain & Note \\
\midrule
Update interval $10\to100$      & $\approx8\times$   & Physics between updates is free \\
Inelastic wall at $0.4A$        & $\approx2\times$   & Adopted from dSBM; one limiter per cell \\
Ramp exponent $\gamma=1\to5$    & $\approx2\times$   & Zero cost; waveform change only \\
Confidence-gated feedback       & $1.45$--$1.94\times$ & Novel; $N\le64$ only \\
\bottomrule
\end{tabular}
\caption{Algorithmic modifications, each measured with the comparator held fixed.}
\label{tab:algo}
\end{table}

\paragraph{Confidence-gated feedback.}
Once a spin's amplitude is pinned at the limiter and its sign has been stable
across consecutive updates, its contribution to every other spin's field is
constant. We compute it once, add it to a cached vector, and remove the spin from
the active set; the matrix multiply then runs over active spins only and shrinks
as the anneal proceeds. This exploits a property specific to analog machines ---
amplitude carries a confidence signal that a purely binary solver does not
possess. Measured MAC reduction is $2.0$--$2.2\times$ with a net TTS improvement
of $1.45$--$1.94\times$ at $N\le64$. The method fails at $N=128$, where baseline
success probability is already $0.01$ and premature freezing destroys the
remaining chance.

Two further variants were tested and closed: parallel tempering across the batch
axis ($0.64$--$0.68\times$, worse) and cyclic pumping ($\le0.13\times$, much
worse). Predictive extrapolation of the feedback was strongly harmful
($p:0.805\to0.105$), as it amplifies measurement noise.

\section{Where an advantage can and cannot exist}
\label{sec:where}

In measurement-feedback form the coupling product $Jx$ is computed digitally, on
an FPGA, exactly as dSBM computes it on a GPU. There is therefore no
per-operation energy advantage available to offset an operation-count deficit,
and Sec.~\ref{sec:parity} shows there is no operation-count surplus either. We
conclude that the measurement-feedback architecture has no defensible energy
claim on dense problems, and we note that energy comparisons in this literature
frequently assume operation parity without testing it.

The remaining possibility is architectural. If the coupling is physically wired
rather than computed --- possible when $J$ is fixed and sparse --- the anneal
requires no digital work at all beyond readout. Under that assumption we compute
a $150$--$180\times$ reduction in digital operations on lattice instances.
Separately, we verified in simulation that such an array anneals successfully
\emph{without} measurement feedback or amplitude error correction, which the
wired configuration cannot support: autonomous operation matched or exceeded
measurement feedback on lattice instances ($p=0.292$ vs.\ $0.254$ at $64$ spins).

We stress that this is an architectural analysis, not a measurement of hardware.

\section{Limitations}
\label{sec:limits}

\begin{itemize}
\item \textbf{All results are simulation.} No hardware was constructed. Every
      statement about physical implementations, settling times, or energy is a
      projection from the model.
\item \textbf{Arm A versus Arm B is untested.} The central hypothesis of this
      architecture class --- that physical analog dynamics outperforms a digital
      simulation of the same dynamics --- cannot be evaluated without hardware
      and is not evaluated here.
\item \textbf{Instance class.} All dense results use random $\pm1$ instances;
      structured, sparse, and application-derived instances may rank algorithms
      differently.
\item \textbf{Comparator implementation.} Our dSBM is a reference
      implementation, not the authors' optimised version, and may understate
      its performance.
\item \textbf{Sizes.} Scaling data ends at $N=64$; extrapolation beyond that is
      not evidence.
\end{itemize}

\section{Conclusion}

The competitive standing of an analog Ising machine against a tuned digital
solver, on dense random instances, is near-parity --- and we could have reported
either a large win or a large loss from the same data by tuning one side. We
believe this is the most useful thing we can contribute: the field should treat
untuned baselines as a defect, state tuning budgets explicitly, and count cost in
operations. We release all code, instances and results so that our numbers can be
checked and our errors found.

\section*{Data and code availability}
All simulation code, the benchmark instance library (100 instances, $N=8$, with
brute-forced ground states, SHA-256
\texttt{b8c6050b2edf1ed40d02ae97e0e37b39ea60661822ccc5afa3c40cfdf3c48812}),
and the complete result set are available at
\url{https://github.com/shmuelbiller1/apim-benchmark-study} and archived at
\url{https://doi.org/10.5281/zenodo.22833988}.

\section*{Acknowledgements}
Portions of the simulation code, analysis and manuscript preparation were carried
out with assistance from a large language model (Anthropic Claude). All
experimental design decisions, result interpretation and conclusions are the
author's responsibility. Several errors in intermediate analyses, including two
of the three benchmark verdicts reported in Table~\ref{tab:invert}, were
identified and corrected during this process; the sequence is retained in the
paper deliberately.

\begin{thebibliography}{9}
\bibitem{goto1959} E.~Goto, ``The parametron, a digital computing element which
utilizes parametric oscillation,'' \emph{Proc.\ IRE}, 1959.
\bibitem{inagaki2016} T.~Inagaki \emph{et al.}, ``A coherent Ising machine for
2000-node optimization problems,'' \emph{Science} \textbf{354}, 603 (2016).
\bibitem{honjo2021} T.~Honjo \emph{et al.}, ``100,000-spin coherent Ising
machine,'' \emph{Science Advances} \textbf{7}, eabh0952 (2021).
\bibitem{leleu2019} T.~Leleu \emph{et al.}, ``Destabilization of local minima in
analog spin systems by correction of amplitude heterogeneity,''
\emph{Phys.\ Rev.\ Lett.} \textbf{122}, 040607 (2019).
\bibitem{goto2019} H.~Goto \emph{et al.}, ``Combinatorial optimization by
simulating adiabatic bifurcations in nonlinear Hamiltonian systems,''
\emph{Science Advances} \textbf{5}, eaav2372 (2019).
\bibitem{goto2021} H.~Goto \emph{et al.}, ``High-performance combinatorial
optimization based on classical mechanics,'' \emph{Science Advances} \textbf{7},
eabe7953 (2021).
\end{thebibliography}

\end{document}
